\documentclass[11pt]{article}

\usepackage[utf8]{inputenc}
\usepackage[T1]{fontenc}
\usepackage{natbib}
\usepackage{geometry}
\geometry{margin=2.5cm}
\usepackage{hyperref}
\usepackage{url}
\usepackage{booktabs}
\usepackage{amsmath}
\usepackage{amsfonts}
\usepackage{amssymb}
\usepackage{microtype}
\usepackage{xcolor}
\usepackage{graphicx}
\usepackage{float}
\usepackage{subcaption}
\usepackage{cleveref}
\crefname{equation}{Eq.\@}{Eqs.\@}
\Crefname{equation}{Eq.\@}{Eqs.\@}
\crefname{figure}{Fig.\@}{Figs.\@}
\Crefname{figure}{Fig.\@}{Figs.\@}
\crefname{table}{Tab.\@}{Tabs.\@}
\Crefname{table}{Tab.\@}{Tabs.\@}
\crefname{section}{Sect.\@}{Sects.\@}
\Crefname{section}{Sect.\@}{Sects.\@}

\def\eg{{\it e.g.}}
\def\ie{{\it i.e.}}
\def\R{\mathbb{R}}
\newcommand{\mlpsub}{\mathrm{MLP}_{\mathrm{sub}}}
\newcommand{\mlpnode}{\mathrm{MLP}_{\mathrm{node}}}

% Figure directory. Change ONLY this line when moving between editors:
%   Overleaf (main.tex next to figures/):  \newcommand{\figdir}{figures/metabolism}
%   local repo (docs/ build):              \newcommand{\figdir}{../figures/metabolism}
\newcommand{\figdir}{../figures/metabolism}

\setlength{\parindent}{0pt}
\setlength{\parskip}{0.5em}

\title{\textbf{Recovering reaction kinetics and homeostatic regulation from
metabolite dynamics with graph neural networks}}
\author{}
\date{}

\begin{document}
\maketitle
\begin{abstract}
A metabolic network is a bipartite graph in which metabolites are coupled
through reactions, each with its own rate constant and stoichiometry. We
pose the \emph{inverse} problem: given only the time series of metabolite
concentrations and the bipartite graph structure, recover the components
that generated the dynamics --- the per-reaction rate constants $k_j$, the
substrate-to-rate kinetic law, and the per-metabolite homeostatic
regulation. We replace the unknown functions of the forward model with
learnable MLPs that operate on the metabolite--reaction graph and a vector
of learnable rate constants, and train the graph neural network (GNN)
end-to-end to predict $dc/dt$. On synthetic mass-action systems with known
ground truth we study rate-constant recovery and the identifiability
obstacles that arise from scale ambiguity, functional compensation, and
homeostasis absorption. We outline the path from synthetic data to real
metabolomics --- starting from fully parameterised kinetic models from
BioModels (yeast glycolysis), constraint-based genome-scale models
(yeast-GEM), and high-frequency real-time metabolomics during nutrient
transitions. This is a companion methods note to our connectome-dynamics
studies, which apply the same sign-/structure-locked GNN inverse-modelling
recipe to neural circuits.
\end{abstract}

\section{Introduction}
\label{sec:intro}

Cellular metabolism is a large, sparse reaction network: hundreds of
metabolites are interconverted by hundreds of enzyme-catalysed reactions,
each proceeding at a rate set by an intrinsic rate constant and the
instantaneous concentrations of its substrates. The \emph{forward} problem
--- predicting concentration dynamics given the network, the kinetics and
the rate constants --- is classical biochemistry. The \emph{inverse}
problem is harder and is our subject: given measured concentration
trajectories and the graph of which metabolites enter which reactions,
recover the rate constants and the functional form of the kinetics. This
is the metabolic analogue of the connectome-to-dynamics inverse problems
we study elsewhere, and it shares their central difficulty: the mapping
from parameters to dynamics is nonlinear, high-dimensional, and degenerate
--- many parameter combinations produce nearly identical trajectories.

\section{Metabolic networks as bipartite graphs}
\label{sec:graph}

A metabolic network is naturally a \textbf{bipartite graph} with two node
types: metabolites and reactions. Edges connect only nodes of different
type, and each edge carries an integer \emph{stoichiometric coefficient}.
A single-partite metabolite--metabolite graph cannot represent the system,
because a reaction couples \emph{several} substrates and products at once
through a shared rate constant: collapsing it to pairwise metabolite edges
would lose the information that those metabolites participate in the same
reaction.

\paragraph{Stoichiometric matrix.}
The connectivity is summarised by the stoichiometric matrix
$\mathbf{S} \in \R^{n \times m}$ ($n$ metabolites, $m$ reactions), where
$S_{ij}$ is the signed coefficient with which metabolite $i$ enters
reaction $j$: negative for substrates (consumed), positive for products
(produced). $\mathbf{S}$ is sparse (each reaction touches only a handful of
metabolites), integer-valued (typically in $\{-2,-1,0,+1,+2\}$), and for
mass-balanced reactions its columns sum to zero.

\section{Forward model: reaction kinetics and homeostasis}
\label{sec:forward}

\paragraph{Mass-action kinetics.}
Under the law of mass action a reaction proceeds faster when its substrates
are more abundant, with each substrate raised to the power of its
stoichiometric coefficient. For reaction $j$ with rate constant $k_j$,
\begin{equation}
v_j \;=\; k_j \prod_{k \in \mathrm{sub}(j)} c_k^{\,|S_{kj}|},
\label{eq:massaction}
\end{equation}
so a reaction consuming two molecules of a substrate depends quadratically
on its concentration, and one consuming a single molecule depends linearly.
The product over substrates (\emph{multiplicative} aggregation) encodes the
requirement that \emph{all} substrates be present simultaneously, and the
resulting nonlinear coupling is what allows autocatalytic cycles to sustain
oscillations. A simpler \emph{additive} aggregation
$v_j = k_j \sum_{k} c_k^{|S_{kj}|}$ treats substrates as contributing
independently and is a useful surrogate for enzyme-mediated reactions that
are not strictly mass-action.

\paragraph{Full dynamics.}
Combining reaction fluxes through $\mathbf{S}$ with a per-metabolite
homeostatic restoring force gives the complete forward model
\begin{equation}
\frac{dc_i}{dt}
  = \underbrace{-\lambda_i\,(c_i - c_i^{\mathrm{baseline}})}_{\text{homeostasis}}
  + \underbrace{\sum_{j=1}^{m} S_{ij}\, v_j}_{\text{reaction dynamics}},
\label{eq:full}
\end{equation}
with $v_j$ from \cref{eq:massaction}. The homeostatic term pulls each
metabolite toward a baseline level $c_i^{\mathrm{baseline}}$ at a
metabolite-specific rate $\lambda_i$; crucially it is \emph{small} relative
to the reaction fluxes, which is what makes it hard to recover
(\cref{sec:identifiability}). \Cref{eq:full} spans a family of regimes used
for synthetic data: pure reaction ($\lambda = 0$, additive), homeostatic
($\lambda > 0$, additive, convergent), and oscillatory (autocatalytic
3-cycles, multiplicative, sustained oscillations).

\section{Inverse problem: the GNN parameterisation}
\label{sec:inverse}

Given observed trajectories $\{c_i(t)\}$ and the bipartite graph, we
recover the generating model by replacing every unknown function in
\cref{eq:full,eq:massaction} with a learnable component on the graph:
\begin{equation}
\frac{dc_i}{dt}
  = \underbrace{\mlpnode(c_i, a_i)}_{\approx\,-\lambda_i(c_i - c_i^{\mathrm{baseline}})}
  + \sum_{j=1}^{m} S_{ij}
    \underbrace{k_j \prod_{k \in \mathrm{sub}(j)} \mlpsub\!\big(c_k, |S_{kj}|\big)}
              _{\approx\,k_j\,c_k^{|S_{kj}|}}.
\label{eq:gnn}
\end{equation}
The GNN operates on the metabolite--reaction graph by message passing:
substrate states are mapped through $\mlpsub$, aggregated per reaction
(multiplicatively for mass-action), scaled by the learnable rate constant
$k_j$, and distributed back to metabolites through $\mathbf{S}$; the
node update $\mlpnode$ adds the self-regulation term. Training minimises the
prediction error on $dc/dt$ end-to-end. The learnable parameters are
listed in \cref{tab:params}.

\begin{table}[H]
\centering
\small
\begin{tabular}{l l l}
\toprule
Parameter & Type & Role \\
\midrule
$a_i \in \R^d$ & embedding vector & per-metabolite identity \\
$k_j$ & scalar ($m$ of them) & per-reaction rate constant (learned in $\log$-space) \\
$\mlpsub(c_k, |S_{kj}|)$ & MLP & discovers the kinetic law $c_k^{|S_{kj}|}$ \\
$\mlpnode(c_i, a_i)$ & MLP & discovers homeostasis $-\lambda_i(c_i - c_i^{\mathrm{baseline}})$ \\
\bottomrule
\end{tabular}
\caption{Learnable components of the GNN inverse model
(\cref{eq:gnn}). The stoichiometric matrix $\mathbf{S}$ is either frozen
from ground truth (``S given'' mode, the focus here) or made learnable
(``S learning'' mode).}
\label{tab:params}
\end{table}

\paragraph{Function discovery.}
The GNN is not told the kinetic law. $\mlpsub$ receives a substrate
concentration and its stoichiometric coefficient and must learn the
power-law response $c_k^{|S_{kj}|}$ from data alone --- in particular it
must distinguish linear ($c^1$) from quadratic ($c^2$) dependence as the
coefficient changes. Likewise $\mlpnode$ must recover the subtle,
small-magnitude homeostatic signal without absorbing dynamics that belong
to the reaction terms.

\section{Identifiability}
\label{sec:identifiability}

The three learnable components interact and can compensate for one another,
producing degenerate solutions with high prediction accuracy but poor
parameter recovery:
\begin{itemize}
  \item \textbf{Scale ambiguity.} The product $k_j\,\mlpsub$ is invariant
    under $k \to \alpha k,\ \mlpsub \to \mlpsub/\alpha$. Without anchoring,
    the MLP absorbs a global scale and shifts every $k_j$. A centring
    regulariser pins $\mathrm{mean}(\log k)$ to the known range centre.
  \item \textbf{Functional compensation.} If $\mlpsub$ settles on a wrong
    exponent (\eg{} $c^{1.5}$ instead of $c^2$), the rate constants can
    partly compensate by re-scaling, masking the error.
  \item \textbf{Homeostasis absorption.} $\mlpnode$ can grow and soak up
    dynamics that should be explained by reactions, biasing the recovered
    rate constants.
\end{itemize}
These are the metabolic counterparts of the sign/scale degeneracies that
motivate the structure-locking and regularisation priors in our
connectome-dynamics models.

\section{Toy model}
\label{sec:toy}

The toy model is a synthetic rank-50 oscillatory mass-action network ($100$
metabolites, $256$ reactions, $\mathbf{S}$ given, no homeostasis), generated with
known rate constants and fit leak-free (the loss reads only the concentration
dynamics, never a ground-truth label). Rate-constant recovery $R^2$ (learned vs.\
true $\log_{10} k_j$) is $0.75 \pm 0.04$ across four seeds (\cref{tab:cv}), and the
two learned functions match ground truth (\cref{fig:toy_dashboard}): $\mathrm{MLP_{sub}}$
recovers $c^{|s|}$ and $\mathrm{MLP_{node}}$ stays $\approx 0$. The free rollout tracks the
oscillations early but, with no multi-step constraint, the autocatalytic dynamics
eventually diverge: this model recovers the parameters but is not a stable simulator.

Does a recurrent curriculum (growing the rollout horizon across epochs, single-step
$\to 1000$ steps, with co-ramped LR and gradient clipping) close the stability gap? A
naive run is a cautionary tale about the metric. The free trajectory no longer diverges
and the \emph{pooled} rollout Pearson even rises, but the pooled number merely rewards
staying bounded; measured \emph{per metabolite}, fidelity \emph{collapses}, and
rate-constant recovery collapses with it (\cref{fig:toy_recurrent}) while
$\mathrm{MLP_{node}}$ drifts to a spurious nonzero function. The multi-step objective admits a
smooth \emph{degenerate} solution --- wrong $k$ plus spurious homeostasis reproduce a damped
trajectory without the mechanism. The lesson: report \emph{per-metabolite} rollout fidelity,
and recognise that the naive curriculum suppresses dynamics rather than buying stability.

\begin{table}[H]
\centering
\small
\begin{tabular}{l c c c}
\toprule
Seed & training completed & rate-constant recovery $R^2$ & rollout Pearson \\
\midrule
77 (published winner) & 100\% & 0.742 & 0.74 \\
42 & 100\% & 0.705 & 0.77 \\
79 & 90\% & 0.808 & 0.37 \\
7  & 90\% & 0.738 & 0.72 \\
\midrule
\multicolumn{2}{l}{\emph{mean $\pm$ s.d.\ (4 mature seeds)}}
  & $\mathbf{0.748 \pm 0.037}$ & $0.65 \pm 0.16$ \\
\midrule
123 & 30\% (stopped) & 0.722 (rising) & 0.65 \\
\bottomrule
\end{tabular}
\caption{Cross-seed reproduction of the leak-free $k_j$-recovery working
point. Rate-constant recovery $R^2$ clusters tightly at $0.75 \pm 0.04$ over
four seeds; seed 123 was stopped early at $0.72$ and still rising. Rollout
Pearson is the global correlation between the learned free rollout and the true
trajectory over the convergent window (before divergence); it is far more
variable ($0.65 \pm 0.16$) and never stable --- every seed recovers the rate
constants but none yields a stable simulator.}
\label{tab:cv}
\end{table}

\begin{figure}[H]
  \centering
  \includegraphics[width=\textwidth]{\figdir/toy_dashboard.png}
  \caption{\textbf{Toy-model results dashboard (leak-free, seed 77).}
  \textbf{(a)} $\mathrm{MLP_{node}}$ per metabolite (homeostasis; dashed: the true zero of
  this regime). \textbf{(b)} $\mathrm{MLP_{sub}}(c,|s|)$ (solid) vs.\ the true power law
  $c^{|s|}$ (dashed) for $|s|{=}1,2$, anchored at $c{=}1$ ($|s|{=}2$ is out-of-support
  extrapolation). \textbf{(c)} Rate-constant recovery: learned vs.\ true $\log_{10} k_j$ for
  all $256$ reactions after the $\mathrm{MLP_{sub}}$ scalar correction (dashed: identity); red
  points exceed the $0.3$-dex outlier gate. \textbf{(d)} Free rollout of the eight
  highest-variance metabolites, $z$-scored and stacked: ground truth (green) vs.\ the learned
  autoregressive rollout (dashed black); red line marks rollout divergence.}
  \label{fig:toy_dashboard}
\end{figure}

\begin{figure}[H]
  \centering
  \includegraphics[width=\textwidth]{\figdir/toy_dashboard_toy_recurrent.png}
  \caption{\textbf{Naive recurrent curriculum: stable but dynamically dead.}
  Same dashboard as \cref{fig:toy_dashboard}, for a model trained with the
  horizon-growing curriculum (single-step $\to 1000$ steps). \textbf{(d)} The free rollout no
  longer diverges but the predicted traces are smooth and bounded, no longer reproducing the
  oscillations; \textbf{(c)} rate-constant recovery has collapsed to a structureless cloud; and
  \textbf{(a)} $\mathrm{MLP_{node}}$ has drifted to a spurious nonzero homeostasis. (\textbf{(b)}
  $|s|{=}2$ is out-of-support extrapolation, as in \cref{fig:toy_dashboard}.)}
  \label{fig:toy_recurrent}
\end{figure}

A different lever closes the stability gap: supplying a known \emph{external stimulus}.
Feeding a time-varying drive to a subset of metabolites at every rollout step --- the
boundary-input convention of the calcium models --- changes the picture. With the drive
supplied and ten driven trajectories for training, even \emph{single-step} training gives a
rollout that is stable over the full window while preserving rate-constant recovery, and stays
bounded on a held-out trajectory with a different drive (\cref{fig:toy_stim}). Running the
stimulus-driven \emph{curriculum} on top (three horizon schedules to a fixed $12$\,h budget;
\cref{fig:toy_stim_rec}) does not improve on this: the stimulus prevents the dead-collapse the
stimulus-free curriculum suffered, but once training reaches a multi-step horizon the
curriculum erodes $k$-recovery without a rollout gain --- the same identifiability tension as
\cref{fig:toy_recurrent}, softened by the drive. \textbf{Single-step training with the
external stimulus remains the best toy configuration.}

\begin{figure}[H]
  \centering
  \includegraphics[width=\textwidth]{\figdir/toy_dashboard_toy_stim10_single.png}
  \caption{\textbf{An external stimulus anchors the rollout.} Same dashboard, for
  single-step training on ten trajectories driven by a known time-varying stimulus (orange in
  \textbf{d}, the given input). \textbf{(d)}~The free rollout (black) is now stable over the
  full window with no divergence (contrast \cref{fig:toy_dashboard}d and
  \cref{fig:toy_recurrent}d); \textbf{(c)}~rate-constant recovery is preserved --- the stimulus
  stabilises the rollout without the parameter collapse the curriculum caused.}
  \label{fig:toy_stim}
\end{figure}

\begin{figure}[H]
  \centering
  \includegraphics[width=\textwidth]{\figdir/toy_dashboard_toy_stim10_cap200.png}
  \caption{\textbf{The recurrent curriculum, even anchored by the stimulus, does not
  beat single-step training.} Best of three stimulus-driven curriculum schedules (capped
  cap-$100$ horizon, $12$\,h budget). \textbf{(d)}~The free rollout (black) under the given
  drive (orange) is stable and tracks the ground truth (green), but \textbf{(c)}~rate-constant
  recovery has degraded (red outliers) and \textbf{(a)}~$\mlpnode$ drifts to a spurious nonzero
  function --- the multi-step objective trades away $k$ identifiability for no rollout gain.}
  \label{fig:toy_stim_rec}
\end{figure}

\textbf{Intrinsic noise as an identifiability lever --- and the right form for a
conserved system.} Stochastic forcing \emph{inside} the dynamics can diversify a
trajectory and help break a degeneracy, but its \emph{form} matters under a conservation
law. Adding noise directly to the concentrations ($dc = f(c)\,dt + \sigma\,dW$) injects mass
no reaction produced --- stoichiometrically unsound. The sound alternative perturbs the
reaction \emph{rates}, $k_j(t) = k_j(1 + \sigma\,\xi_j(t))$ (enzyme-activity fluctuation): mass
is conserved through $S$, trajectories stay physical, and the data diversify
(\cref{fig:toy_knoise}). But here this does not change the outcome: with $S$ given, $k$ is
\emph{already} identifiable at $\sigma = 0$, so there is no degeneracy left to break, and the
noise is benign at low level then harmful as it corrupts the trajectory. The lever would
matter only in a genuinely $k$-degenerate regime (e.g.\ $S$ unknown), left to future work.

\begin{figure}[H]
  \centering
  \includegraphics[width=\textwidth]{\figdir/toy_knoise_rollout.png}
  \caption{\textbf{Stoichiometrically-sound rate-noise sweep (rollout traces).}
  Intrinsic noise applied to the reaction \emph{rates}, $k_j(1+\sigma_k\xi)$
  (enzyme-activity fluctuation; mass-conserving), at four levels. Each panel shows the noisy
  simulated ground truth (green) and the trained model's deterministic free rollout (black),
  $z$-scored and stacked, with $k$-recovery $R^2$ and per-metabolite rollout Pearson annotated.
  As $\sigma_k$ grows the simulation becomes visibly jagged while the model rollout stays
  smooth.}
  \label{fig:toy_knoise}
\end{figure}

\section{Toward real metabolomics}
\label{sec:realdata}

\subsection{First look: are the four real datasets amenable?}
\label{sec:amenability}

Before training, we run the same cheap four-probe test on \emph{all four} datasets measured by
\citet{link2015} --- \emph{E.\ coli}, \emph{B.\ subtilis}, \emph{Y.\ lipolytica}, and a hybrid
sample --- asking whether each recording carries the structure the GNN needs
(\cref{tab:amenability}, \cref{fig:amenability_grid}).

\paragraph{What the three probes compute, and how to read them.}
\begin{itemize}
  \item \textbf{SVD low-rank} ($\mathrm{rank}_{99}$, the number of modes for $99\%$ of the
    $z$-scored variance): the effective dimensionality of the dynamics. A small
    $\mathrm{rank}_{99}/N$ means coordinated, low-dimensional motion a compact network model can
    fit; $\approx N$ means no parsimonious mechanism exists.
  \item \textbf{Per-ion lag-1 autocorrelation} $\rho_1$ (median over ions): temporal smoothness.
    $\rho_1\to1$ is a smooth trace whose derivative $dc/dt$ carries signal; $\rho_1\to0$ is
    white-noise-like, so $dc/dt$ is noise. This is the decisive test, since the GNN learns
    $dc/dt=f(c)$.
  \item \textbf{KEGG coverage}: the fraction of ions mapping (by $m/z$ annotation) onto a
    reference network (iJO1366) --- whether a stoichiometry $\mathbf{S}$ can be built for the
    measured species at all.
\end{itemize}

\begin{table}[H]
\centering
\small
\begin{tabular}{l c c c c c}
\toprule
Organism & ions$\times$t & $\mathrm{rank}_{99}$ & autocorr med. & in iJO1366 & $dc/dt$ \\
\midrule
\emph{E.\ coli}       & $247{\times}119$ & 102 (41\%) & 0.27 & 190 (77\%) & noisy \\
\emph{B.\ subtilis}   & $239{\times}119$ & 98 (41\%)  & 0.41 & 160 (67\%) & noisy \\
\emph{Y.\ lipolytica} & $164{\times}59$  & 51 (31\%)  & \textbf{0.62} & 103 (63\%) & \textbf{smooth} \\
Hybrid                & $338{\times}60$  & 55 (16\%)  & 0.24 & 155 (46\%) & noisy \\
\bottomrule
\end{tabular}
\caption{Amenability of the four measured datasets. $\mathrm{rank}_{99}$ is the
number (and fraction) of singular components for $99\%$ variance; autocorr is
the median per-ion lag-1 autocorrelation; coverage is the count mapped onto
iJO1366. All four are low-rank and reasonably network-mappable; only
\emph{Y.\ lipolytica} has smooth ($dc/dt$-learnable) dynamics.}
\label{tab:amenability}
\end{table}

\noindent\textbf{Does it generalise?} Two signals hold across all four datasets --- the dynamics
are low-rank and a large share of metabolites are core/conserved --- but the third, temporal
smoothness of $dc/dt$, does \emph{not}: \emph{E.\ coli}, \emph{B.\ subtilis} and the hybrid are
noise-dominated, whereas \textbf{\emph{Y.\ lipolytica} is smooth}. The actionable consequence: the
most promising real-data target is not \emph{E.\ coli} (best coverage but noisiest) but
\emph{Y.\ lipolytica}, restricted to its smooth, network-mapped ions and judged by held-out
prediction.

\begin{figure}[H]
  \centering
  \includegraphics[width=0.49\textwidth]{\figdir/amenability_ecoli.png}\hfill
  \includegraphics[width=0.49\textwidth]{\figdir/amenability_bsubt.png}\\[2pt]
  \includegraphics[width=0.49\textwidth]{\figdir/amenability_ylipo.png}\hfill
  \includegraphics[width=0.49\textwidth]{\figdir/amenability_hybr.png}
  \caption{\textbf{Identical amenability probes across the four measured
  datasets:} \textbf{(a)} \emph{E.\ coli}, \textbf{(b)} \emph{B.\ subtilis}, \textbf{(c)}
  \emph{Y.\ lipolytica}, \textbf{(d)} hybrid. Within each panel, top-left: SVD cumulative
  variance; top-right: per-ion lag-1 autocorrelation; bottom-left: six highest-variance traces
  ($z$-scored); bottom-right: KEGG coverage against iJO1366.}
  \label{fig:amenability_grid}
\end{figure}


The framework is validated on synthetic data with known ground truth; the goal is real
metabolomic time series. Three vendored datasets stage that transition from fully parameterised
kinetics, through genome-scale topology, to a real recording. None is solved yet; this section
describes what each provides and what stands between it and a result.

\paragraph{Rung 1 --- yeast glycolysis kinetic model (BioModels
the published yeast glycolysis kinetic model).}
A small, fully parameterised kinetic model of glycolysis in \emph{S.\ cerevisiae}
\citep{smallbone2013}, imported from SBML: $20$ internal metabolites ($+9$ fixed boundary
species), $30$ reactions including iso-enzymes, with real stoichiometry $\mathbf{S}$ and
measured initial concentrations. Unlike the all-$|s|{=}1$ toy it is genuinely
\emph{multi-substrate} ($15$ first-order, $13$ second-order, one fourth-order), and the rate
laws are \emph{Michaelis--Menten} (MM), the saturating law $v = V_{\max}\,c/(K_m+c)$, so
$v_j = V_{\max,j}\prod_k[c_k/(K_{m,kj}+c_k)]^{|S_{kj}|}$ rather than a power law. $V_{\max}$ is
the per-reaction rate constant to recover (the MM analogue of $k_j$) and $K_m$ the
half-saturation constant; known $V_{\max}/K_m$ make this the cleanest available ground truth.
\emph{Status:} the SBML import, $\mathbf{S}$ extraction, and an MM simulator are in place
($12{,}000$-frame dataset), with the connectome-style autoregressive curriculum ported to
stabilise the rollout. Two cautions, both found by being sceptical of a good number. A test
rollout that looked catastrophic was an evaluator bug (the boundary stimulus was omitted);
supplied, the curriculum model reproduces the trajectory well (\cref{fig:glyco_rollout}). But
that success is \emph{in-sample}, and both the leak-resistant recovery test
(\cref{fig:glyco_recovery}) and a direct held-out rollout (\cref{fig:glyco_holdout}) fail: the
model reproduces the dynamics with the \emph{wrong} kinetics, a degenerate fit, and the good
in-sample rollout was effectively memorisation of one trajectory. The curriculum fixed rollout
\emph{stability}; the open problem on real MM data is parameter \emph{identifiability} --- the
same scale/compensation degeneracy as the synthetic setup, here unbroken.

\begin{figure}[H]
  \centering
  \includegraphics[width=0.49\textwidth]{\figdir/glyco_rollout.png}\hfill
  \includegraphics[width=0.49\textwidth]{\figdir/glyco_rollout_holdout.png}
  \caption{\textbf{Rung 1 glycolysis rollout: in-sample vs held-out.} Eight highest-variance
  metabolites, $z$-scored and stacked: ground truth (green) vs the learned free autoregressive
  rollout (dashed black); the given external stimulus channels (orange, top) are the known
  boundary drive, not predicted. \textbf{Left (in-sample):} the rollout tracks the training
  trajectory. \textbf{Right (held-out):} on an unseen trajectory (different stimulus and
  initial state) the rollout settles to the wrong levels and no longer tracks ground truth.}
  \label{fig:glyco_rollout}
  \label{fig:glyco_holdout}
\end{figure}

\begin{figure}[H]
  \centering
  \includegraphics[width=0.6\textwidth]{\figdir/glyco_k_recovery.png}
  \caption{\textbf{Rung 1 glycolysis: $V_{max}$ recovery (leak-resistant test).} Learned vs
  true $\log_{10} V_{max}$ for the $30$ reactions (dashed: identity); red points exceed the
  $0.3$-dex outlier gate. Despite the near-perfect in-sample rollout (\cref{fig:glyco_rollout}),
  the recovered kinetics do not track ground truth --- a degenerate fit.}
  \label{fig:glyco_recovery}
\end{figure}

\paragraph{Does $\mathrm{MLP_{sub}}$ learn the glycolysis kinetics, and does the
log-exp form help?} A natural suspect for the recovery failure is that $\mathrm{MLP_{sub}}$ cannot represent the
saturating MM substrate law $[c/(K_m{+}c)]^{|s|}$. Comparing the learned $\mathrm{MLP_{sub}}$
against that truth, for the plain MLP and the log-space form $g=\exp(\mathrm{MLP}(\log c, |s|))$
(\cref{fig:glyco_sub}), representation is \emph{not} the bottleneck: the well-populated $|s|{=}1$
and $|s|{=}2$ shapes are captured, and only the single $|s|{=}4$ reaction is unconstrained (too
few examples to fit). The log-exp form does not rescue recovery either. The limiting factor is
not how $\mathrm{MLP_{sub}}$ is parameterised but the MM identifiability degeneracy ($V_{\max}$ and
$K_m$ trade off): the substrate function can be represented, yet the rate constants are not
recovered.

\begin{figure}[H]
  \centering
  \includegraphics[width=\textwidth]{\figdir/glyco_sub_compare.png}
  \caption{\textbf{Rung 1: does $\mathrm{MLP_{sub}}$ learn the glycolysis substrate
  kinetics?} Learned $\mathrm{MLP_{sub}}(c,|s|)$ (solid) vs the true Michaelis--Menten law
  $[c/(K_m{+}c)]^{|s|}$ (dashed) for the substrate orders $|s|{=}1,2,4$, log $y$-axis, anchored
  at $c{=}1$. \textbf{(a)} plain MLP, \textbf{(b)} log-space MLP. The well-populated $|s|{=}1,2$
  shapes are captured; the single $|s|{=}4$ reaction is unconstrained.}
  \label{fig:glyco_sub}
\end{figure}

\paragraph{The limit is signal, not capacity.} Two controls on a synthetic $|s|{\ge}2$ regime localise the residual failure. Training on four
diverse trajectories instead of one barely helps, so it is not a data-volume problem; and a
\emph{structured} substrate function $g=c^{a(|s|)}$ --- which can represent any exponent exactly
--- still fails, its learned exponent map nailing $a(1)$ but never reaching the high exponents
because those reactions are \emph{rare}. The bottleneck for high-order kinetics is neither MLP
capacity nor data volume but the \emph{scarcity of high-order reactions}: too little signal to
identify a steep exponent. The fix is a structural prior (e.g.\ integer exponents) or a regime
richer in high-$|s|$ reactions, not a better function approximator.

\paragraph{Rung 2 --- consensus genome-scale model (yeast-GEM).}
The community-curated genome-scale model of \emph{S.\ cerevisiae} \citep{yeastgem}: $2{,}806$
metabolites, $4{,}131$ reactions, supplying a large real $\mathbf{S}$ and flux bounds but
\emph{no} kinetic rate laws. The intended use is to extract a central-carbon subgraph, drive it
with dynamic-FBA or literature kinetics, and fit the GNN in ``S given'' mode on realistic
topology. \emph{Status:} not started beyond vendoring; subgraph extraction and a dynamics
generator remain to be written.

\paragraph{Rung 3 --- real-time metabolomics during nutrient transitions
(real data).}
High-frequency flow-injection mass spectrometry of live microbial culture \citep{link2015}.
For \emph{E.\ coli} the released data give $247$ annotated ions over $119$ time points in three
replicates, plus a starvation$\to$growth switch series; companion \emph{B.\ subtilis} and
\emph{Y.\ lipolytica} sets are included. Each ion carries a KEGG annotation, though many are
mass-ambiguous and the values are \emph{relative} intensities, not absolute concentrations.
\emph{Status:} ingested and characterised (\cref{fig:ecoli_real}). The trajectory is
\emph{information-rich} (activity rank well above the closed synthetic networks), so data
richness is \emph{not} the limiting factor. The obstacle is \emph{partial observation}: the
measured set is a small, open fraction of iJO1366 whose dynamics depend on unmeasured species;
the only fully-observed subnetwork (reactions whose every metabolite is measured) is $\sim181$
of iJO1366's reactions, the immediate LSQ/GNN target. There is \emph{no} kinetic ground truth,
so success can be judged only by held-out prediction. A parallel mammalian direction is chronic
jugular microdialysis in freely moving mice \citep{nardin2025}.

\begin{figure}[H]
  \centering
  \includegraphics[width=\textwidth]{\figdir/ecoli_real_data.png}
  \caption{\textbf{The real \emph{E.\ coli} metabolome is a high-rank trajectory.}
  \textbf{(a)} Measured trajectory: $247$ metabolites $\times$ $119$ time points (relative ion
  intensities; $z$-scored per metabolite, rows sorted by activity). \textbf{(b)} SVD
  cumulative-variance spectrum (activity rank far richer than the closed synthetic networks).}
  \label{fig:ecoli_real}
\end{figure}

\paragraph{What transfers and what changes.}
The core recipe transfers: the bipartite graph, message passing over reactions, function
discovery via $\mlpsub/\mlpnode$, and log-space rate-constant recovery. Real data adds, all at
once, the obstacles the synthetic setup omits: saturating/allosteric kinetics that $\mlpsub$
must learn; a partial, incomplete $\mathbf{S}$; far fewer, unevenly spaced time points; relative
rather than absolute concentrations; and external drivers entering as $f_{\mathrm{ext}}(t)$ in
\cref{eq:full}.

\paragraph{Identifiability on a real network topology (imposed kinetics).}
Short of full real data, we can probe identifiability on \emph{real} topology by imposing
synthetic MM kinetics (random $V_{\max}/K_m$, hence known ground truth) on the real $\mathbf{S}$
of two networks: \emph{E.\ coli} core ($72$ metabolites, $71$ reactions) and a yeast-GEM
central-carbon subgraph ($208$ metabolites, $120$ reactions). With $\mathbf{S}$ given the model
reproduces the dynamics yet recovers \emph{none} of the kinetics on either network
(\cref{fig:rung_real}). The cause is dynamical, not structural: these closed networks relax to a
steady state within a few frames, so their trajectories have activity rank ${\approx}2$ --- far
below the oscillatory toy. Identifiability tracks the \emph{information content of the dynamics},
not network size or method: a closed subnetwork at steady state is information-starved.
Recovering kinetics on real topology requires holding the system \emph{off} its fixed point ---
a sustained external drive, or the nutrient switch of Rung~3 --- to raise the rank.

\begin{figure}[H]
  \centering
  \includegraphics[width=0.49\textwidth]{\figdir/toy_dashboard_ecoli_core_mm.png}\hfill
  \includegraphics[width=0.49\textwidth]{\figdir/toy_dashboard_yeast_central_mm.png}
  \caption{\textbf{Identifiability on real topology with imposed MM kinetics.} Left:
  \emph{E.\ coli} core ($72{\times}71$); right: yeast-GEM cytosolic subgraph ($208{\times}120$).
  Real stoichiometry, synthetic MM kinetics (so $V_{\max}/K_m$ are known), $\mathbf{S}$ given.
  Each dashboard: \textbf{(a)} $\mlpnode$, \textbf{(b)} $\mlpsub$, \textbf{(c)} $V_{\max}$
  recovery, \textbf{(d)} free rollout.}
  \label{fig:rung_real}
\end{figure}

\paragraph{Raising identifiability: drive the network off steady state.}
The failure above is a \emph{data} problem with a precise linear-algebra reading. With
$\mathbf{S}$ given, recovering the kinetics is \emph{linear in the parameters},
$\mathbf{A}\,\boldsymbol{\theta}=\mathbf{b}$; a metabolite that never varies gives constant or
zero columns --- null directions in $\mathbf{A}$ --- so its parameters are unidentifiable.
Equivalently, linearising as $\dot c = J c + B\,u(t)$, an autonomous trajectory collapses into
a few slow eigenmodes, whereas an external source $u$ on $m$ metabolites drives it into a
controllability subspace whose dimension grows with $m$. Injecting a known multi-frequency
metabolite source and sweeping its breadth, the activity rank rises \emph{linearly} with the
number of driven metabolites (rank ${\approx}m$) on both real topologies (\cref{fig:stim_rank}),
clearing the rank-$20$ target. Information content, not network size, is the lever; the next
step trains on these high-rank, stimulus-driven trajectories.

\begin{figure}[H]
  \centering
  \includegraphics[width=0.62\textwidth]{\figdir/stimulus_rank.png}
  \caption{\textbf{An external metabolite source raises the activity rank
  ${\approx}$ linearly.} Activity rank ($99\%$) of the simulated trajectory vs the number of
  metabolites carrying a known multi-frequency source drive, for all three rungs under imposed
  MM kinetics: Rung~1 glycolysis ($20$ met), Rung~3 \emph{E.\ coli} core ($72$ met), Rung~2
  yeast-GEM subgraph ($208$ met). Dashed line: rank ${\approx}m$ identity.}
  \label{fig:stim_rank}
\end{figure}

\begin{figure}[H]
  \centering
  \includegraphics[width=\textwidth]{\figdir/model_traces.png}
  \caption{\textbf{Free-rollout traces under the external drive, one panel per rung} (the
  trained GNN). Eight highest-variance metabolites, $z$-scored and stacked: ground truth (green)
  vs the model's free rollout (black). \textbf{(a)} Rung~1 glycolysis; \textbf{(b)} Rung~3
  \emph{E.\ coli} core and \textbf{(c)} Rung~2 yeast-GEM subgraph, both with the $m{=}40$
  metabolite source holding the network in sustained oscillation. Per-metabolite Pearson
  annotated. The least-squares companion is \cref{fig:lsq_rollout}.}
  \label{fig:model_traces}
\end{figure}

\paragraph{Identifiability is not the bottleneck --- learnability is.} The GNN trained on the
driven (rank-$38$) \emph{E.\ coli} data still recovers \emph{none} of the kinetics, yet the
inverse problem is \emph{well-posed}: given $\mathbf{S}$ \emph{and the rate-law shape}, $dc/dt$ is
linear in the per-reaction scale, $A\,\theta = b$ with $A_{\cdot j} = \mathbf{S}_{\cdot j}\,\phi_j(c)$,
and built from the simulator's Jacobian this $A$ is full column rank and well conditioned at both
rank-$2$ and rank-$38$ --- so a plain least-squares solve recovers $V_{\max}$ exactly and rolls it
out perfectly (\cref{fig:lsq_rollout}). The catch is that this is conditional on \emph{knowing
$\phi$}; the GNN must learn the shape ($\mathrm{MLP_{sub}}$) and the scale ($V_{\max}$) jointly, and
a saturating shape with a free $K_m$ can absorb the scale so $V_{\max}$ drifts. GNN success is thus
set by shape$\times$scale learnability, not identifiability. The lever is the model, not the data:
solve $V_{\max}$ by least squares given the learned shape, or constrain $\mathrm{MLP_{sub}}$ --- not
more rank, drive, or trajectories.

\begin{figure}[H]
  \centering
  \includegraphics[width=\textwidth]{\figdir/lsq_rollout_traces.png}
  \caption{\textbf{Least-squares rollout (companion to \cref{fig:model_traces}).} Free-rollout
  traces for the same three rungs and external drive as \cref{fig:model_traces}, but with the
  per-reaction $V_{\max}$ set by a least-squares solve given the true substrate shape rather than
  by the GNN. Eight highest-variance metabolites, $z$-scored and stacked: ground truth (green) vs
  the least-squares rollout (black). \textbf{(a)} Rung~1 glycolysis; \textbf{(b)} Rung~3
  \emph{E.\ coli} core; \textbf{(c)} Rung~2 yeast-GEM subgraph. Per-metabolite Pearson annotated.}
  \label{fig:lsq_rollout}
\end{figure}

\paragraph{A realistic, non-oscillatory stimulus changes nothing.} One might worry the failure is
an artefact of the \emph{shape} of the drive --- a clean multi-sinusoid is nothing like a real
metabolome. Replacing the sinusoid with an Ornstein--Uhlenbeck / AR($1$) drive ($\phi{=}0.98$),
smooth and aperiodic like the measured \emph{E.\ coli} trajectory, and regenerating all three rungs
with the same imposed-MM kinetics reproduces the failure \emph{exactly}: the GNN still recovers
$V_{\max}$ at $R^2 \approx 0$ while rolling out serviceably. Least squares, by contrast, recovers and
free-rolls-out the dynamics almost perfectly on the same realistic-stimulus data
(\cref{fig:lsq_rollout_ou}), exactly as under the sinusoid (\cref{fig:lsq_rollout}). The learnability
barrier is \emph{stimulus-shape-invariant}: neither the richness nor the temporal character of the
drive is the lever. The lever remains model-side.

\begin{figure}[H]
  \centering
  \includegraphics[width=\textwidth]{\figdir/lsq_rollout_traces_ou.png}
  \caption{\textbf{Least-squares rollout under a realistic (OU) stimulus} (companion to
  \cref{fig:lsq_rollout}). As \cref{fig:lsq_rollout}, but on the realistic Ornstein--Uhlenbeck /
  AR($1$)-driven rungs: ground truth (green) vs the least-squares rollout (black), $z$-scored and
  stacked. \textbf{(a)} Rung~1 glycolysis; \textbf{(b)} Rung~2 \emph{E.\ coli} core; \textbf{(c)}
  Rung~3 yeast central. Per-metabolite Pearson annotated.}
  \label{fig:lsq_rollout_ou}
\end{figure}

\paragraph{Least squares does the job; the GNN hardly does it yet.} Putting the two estimators
side by side on the same realistic-stimulus data (\cref{fig:param_recovery_lsq,fig:param_recovery_gnn})
locates the difficulty. Given the true substrate shape, least squares recovers $V_{\max}$ exactly on
the two rungs where the dynamics are linear in $V_{\max}$ (glycolysis and \emph{E.\ coli} core); the
third (yeast central) is flux-limit-clamped, a simulator artefact, not a data limit. The GNN, which
must \emph{discover} shape and scale jointly, recovers essentially none of it. The remaining gap is
closed model-side: feed the GNN the shape (or constrain $\mathrm{MLP_{sub}}$) and solve the scale by
least squares, rather than asking it to learn both at once.

\begin{figure}[H]
  \centering
  \includegraphics[width=\textwidth]{\figdir/param_recovery_lsq.png}
  \caption{\textbf{Per-reaction $V_{\max}$ recovery --- least squares.} Recovered vs true
  $\log_{10} V_{\max}$ for every reaction (least squares given the true substrate shape), one panel
  per rung: \textbf{(a)} glycolysis, \textbf{(b)} \emph{E.\ coli} core, \textbf{(c)} yeast central.
  Dashed line: exact recovery; black inliers, red points exceed the $0.3$-dex outlier gate (raw and
  trimmed $R^2$ and outlier fraction annotated). Companion to \cref{fig:param_recovery_gnn}.}
  \label{fig:param_recovery_lsq}
\end{figure}

\begin{figure}[H]
  \centering
  \includegraphics[width=\textwidth]{\figdir/param_recovery_gnn.png}
  \caption{\textbf{Per-reaction $V_{\max}$ recovery --- the trained GNN.} As
  \cref{fig:param_recovery_lsq}, for the GNN's leak-resistant recovery: recovered vs true
  $\log_{10} V_{\max}$ per reaction, one panel per rung. Dashed line: exact recovery.}
  \label{fig:param_recovery_gnn}
\end{figure}

\begin{table}[H]
\centering
\small
\resizebox{\textwidth}{!}{%
\begin{tabular}{l l l l}
\toprule
Target & Network / kinetics & Ground truth & Status \\
\midrule
Synthetic $k_j$ (S given) & random / mass-action & $k_j$ known & \textbf{works}, $R^2{\approx}0.87$ \\
Synthetic homeostasis & random / mass-action & $\lambda_i$ known & open (single-step blind) \\
Rung 1: yeast glycolysis & real / MM+Hill & $V_{\max},K_m$ known & rollout in-sample; $V_{\max}$ R$^2{\approx}0$ (degenerate) \\
Real topology (\emph{E.\ coli} core) & real S / imposed MM & $V_{\max}$ known & $V_{\max}$ R$^2{\approx}0$, $100\%$ out (rank-2 degenerate) \\
Real topology (yeast-GEM subgraph) & real S / imposed MM & $V_{\max}$ known & $V_{\max}$ R$^2{\approx}0$, $100\%$ out (rank-2 degenerate) \\
\quad + realistic OU drive (rungs 1--3) & real S / imposed MM & $V_{\max}$ known & $V_{\max}$ R$^2{\approx}0$; rollout $0.83$--$0.90$ (shape-invariant) \\
Rung 3: real \emph{E.\ coli} data & inferred / unknown & none & ingested; rank $102$, $181$ fully-obs.\ rxns; LSQ/GNN next \\
\bottomrule
\end{tabular}}
\caption{Where each target stands. Only the first row is a validated,
leak-free result; the rest is the work ahead.}
\label{tab:status}
\end{table}


\begin{thebibliography}{9}

\bibitem[Smallbone et al.(2013)]{smallbone2013}
Smallbone, K., Messiha, H.L., Carroll, K.M., \emph{et al.} (2013).
A model of yeast glycolysis based on a consistent kinetic characterisation
of all its enzymes. \emph{FEBS Letters} 587(17), 2832--2841.
BioModels the published yeast glycolysis kinetic model.

\bibitem[Lu et al.(2019)]{yeastgem}
Lu, H., Li, F., S{\'a}nchez, B.J., \emph{et al.} (2019).
A consensus S.\ cerevisiae metabolic model Yeast8 and its ecosystem for
comprehensively probing cellular metabolism. \emph{Nature Communications}
10, 3586. yeast-GEM, SysBioChalmers.

\bibitem[Link et al.(2015)]{link2015}
Link, H., Fuhrer, T., Gerosa, L., Zamboni, N., \& Sauer, U. (2015).
Real-time metabolome profiling of the metabolic switch between starvation
and growth. \emph{Nature Methods} 12, 1091--1097.
\url{https://www.nature.com/articles/nmeth.3584}.

\bibitem[Nardin et al.(2025)]{nardin2025}
Nardin, M., Wang, N., Elziny, S., \emph{et al.} (2025).
Reconstructing a physiological state space via chronic jugular
microdialysis in freely moving mice. \emph{bioRxiv}.
\url{https://doi.org/10.64898/2025.12.08.692974}.

\end{thebibliography}

\end{document}
